\chapter{随机过程}
\section{马尔可夫链}
\noindent\textbf{要点与方法.}\enspace
马尔可夫建模先确定哪些历史信息影响下一步：模式等待保留有效后缀，递归任务保留类型，维修系统还须计入跳转前的等待时间。写出首步方程和吸收边界后，求矩之前可用固定步数内统一的结束概率证明几何尾。

平稳分布须满足平衡方程与归一化。有限不可约链的长期占用比例由它给出，但分布收敛还须排除周期等障碍；清空、覆盖旧记录等结构可用于耦合。固定窗口计数的方差要计入协方差，多目标覆盖则可结合容斥与尾和。
\par\smallskip



\problem{R01}{滑动窗口报警}{状态压缩}
独立试验成功概率为 $p\in(0,1)$，失败概率为 $q=1-p$。从空记录开始，只要最近至多三次试验中已有两次成功便报警；因此 HH 或 HTH 都能触发。求报警时间的均值，并求最后由 HH 触发的概率。
\solution 未报警时只需记录三种状态：0表示最近两次没有 H，A表示刚出现 H，B表示末尾为 HT；空记录也属于0。令剩余时间均值为 $e_0,e_A,e_B$。首步分解给出
\[
 e_0=1+qe_0+pe_A,\quad e_A=1+qe_B,\quad e_B=1+qe_0.
\]
连续两次 H 必能结束，故等待时间有几何尾，上述均值有限。代入后
\[
 (1-q^2)e_0=\frac1p+1+q,\qquad
 \boxed{e_0=\frac{1/p+2-p}{p(2-p)}}.
\]
设从各状态最终由 HH 触发的概率为 $u_0,u_A,u_B$。则 $u_0=u_A$、$u_A=p+qu_B$、$u_B=qu_0$，故
\[
 \boxed{u_0=\frac{p}{1-q^2}=\frac1{2-p}}.
\]
公平情形的平均时间为 $14/3$，HH 触发概率为 $2/3$。
\insight 需要保留一次失败后的有效信息；只记连续成功次数会把 HT 与无有效前缀的状态混淆。

\problem{P16}{同时收集两个重复模式}{模式覆盖}
不断独立抛掷公平硬币，首次在记录中既出现过 HH 又出现过 TT 时停止。两个模式可以重叠，记录不清空。求总抛掷次数 $T$ 的均值与方差。
\solution 先等第一次出现相邻相同。第一枚之后，每次与前一枚相同的指示变量是独立的公平伯努利变量：给定第一枚和全部切换指示，整串硬币唯一确定，且每串等概率。因此该阶段长度 $D=1+G$，其中 $G$ 在正整数上服从参数 $1/2$ 的几何分布，故 $\E D=3,\Var D=2$。

第一次出现 HH 后，还须等 TT；首次出现 TT 后同理。设 $U$ 为第二阶段长度。以“末尾没有目标字母”和“末尾已有一个目标字母”为0、1，均值满足
\[
 a_0=1+\tfrac12a_0+\tfrac12a_1,\qquad
 a_1=1+\tfrac12a_0,
\]
故 $a_0=6,a_1=4$。连续两个目标字母提供几何尾，所有矩有限。二阶矩 $b_0,b_1$ 满足
\[
 b_0=1+a_0+a_1+\tfrac12b_0+\tfrac12b_1,
 \quad b_1=1+a_0+\tfrac12b_0.
\]
代入得 $b_0=58$，所以 $\Var U=58-36=22$。第一阶段结束后未来硬币独立于过去；由 H、T 对称性，第二阶段的条件分布不依赖首个模式或 $D$，故 $D,U$ 独立。于是
\[
 \boxed{\E T=9},\qquad \boxed{\Var T=24}.
\]
\insight “两个模式都出现”是覆盖条件，不能把两个从空记录开始的等待时间直接相加；第二阶段的起始后缀由第一阶段决定。


\problem{R02}{随机清空队列的平稳分布}{重置与耦合}
队列长度 $X_n\in\{0,1,\ldots\}$ 每步独立选择操作：以概率 $a$ 增加1，以概率 $b$ 减少1（零时不动），以概率 $c$ 清空。设 $a,b,c>0$ 且和为1。求平稳分布、平稳平均长度，并给出从任意初态到平稳分布的全变差收敛界。
\solution 对 $k\ge1$，平稳方程为 $\pi_k=a\pi_{k-1}+b\pi_{k+1}$。令
\[
 r=\frac{1-\sqrt{1-4ab}}{2b}\in(0,1),
 \qquad \boxed{\pi_k=(1-r)r^k}.
\]
这里 $br^2-r+a=0$；该多项式在0为正、在1为 $-c$，故所取根位于 $(0,1)$。代入可验证正状态方程。零状态的方程为 $\pi_0=b(\pi_0+\pi_1)+c$，也由同一根等式成立；几何级数保证总和为1。因此
\[
 \boxed{\E_\pi X=\frac r{1-r}}.
\]
将任意初态链与平稳初态链使用同一串操作。第一次清空后，两条链永远相同。前 $n$ 步没有清空的概率为 $(1-c)^n$，由耦合不等式，
\[
 \boxed{\|\mathcal L(X_n)-\pi\|_{\rm TV}\le(1-c)^n}.
\]
这同时证明平稳分布唯一。链不可约且正再生，故从0开始的正时间首次返回均值为 $1/\pi_0=1/(1-r)$；一次清空或空队列服务也算返回。
\insight 重置提供了与初始队长无关的遗忘机制；稳态尾部和收敛速度是两个不同的量。

\problem{R03}{两类递归任务的总工作量}{多类型递推}
一个 A 类任务处理后，以各 $1/2$ 的概率不再产生任务或产生一个 A 和一个 B。一个 B 类任务以概率 $2/3$ 不再产生任务，以概率 $1/3$ 产生两个 A。各任务的选择独立，处理每个任务计一次操作。从一个 A 开始，证明所有任务几乎必然处理完，求两类任务的期望总数；若 A、B 的单位成本为 $c_A,c_B\ge0$，求期望总成本。
\solution 先给未完成的 A、B 任务分别赋权9、7。处理 A 后子任务的期望权重为 $(9+7)/2=8$，处理 B 后为 $2\cdot9/3=6$，均比被处理任务少1。设全部完成时刻为 $T$，在前 $n\wedge T$ 次处理后剩余权重非负，逐步取条件期望得
\[
 \E(n\wedge T)\le9.
\]
单调收敛给出 $\E T\le9$，从而 $T<\infty$ 几乎必然成立。于是各类总数均可积。记其均值为 $A,B$。按所有任务产生的子任务计数，
\[
 A=1+\tfrac12A+\tfrac23B,\qquad B=\tfrac12A.
\]
代入得 $A=1+5A/6$，因此
\[
 \boxed{A=6,\quad B=3,\quad \E T=9},\qquad
 \boxed{\E C=6c_A+3c_B}.
\]
\insight 类型会影响后续分枝，不能只用“平均产生几个子任务”而忽略类型；正权重先证明可积性，再解计数方程。

\problem{R04}{随机翻转位的双边吸收}{生灭链}
有 $m\ge2$ 个二进制位，每步均匀选一位翻转。初始有 $k$ 个1，其中 $1\le k<m$；首次全为0或全为1时停止。求先全为1的概率 $h_k$。另求 $m=4,k=2$ 时的平均停止步数。
\solution 只记录1的个数即可。任意未吸收状态都有概率至少 $m^{-m}$ 在随后至多 $m$ 步到达全零，因此停止时间有几何尾。对 $1\le k<m$，
\[
 h_k=\frac{k}{m}h_{k-1}+\frac{m-k}{m}h_{k+1},
 \qquad h_0=0,\quad h_m=1.
\]
令 $d_k=h_k-h_{k-1}$。则 $(m-k)d_{k+1}=kd_k$，所以
\[
 d_k=\frac{d_1}{\binom{m-1}{k-1}},\qquad
 \boxed{h_k=\frac{\displaystyle\sum_{j=0}^{k-1}\binom{m-1}{j}^{-1}}
 {\displaystyle\sum_{j=0}^{m-1}\binom{m-1}{j}^{-1}}}.
\]
当 $m=4$，对称性给出剩余时间 $e_1=e_3$，且
\[
 e_1=1+\tfrac34e_2,\qquad e_2=1+e_1.
\]
故 $e_1=7$，$\boxed{e_2=8}$。吸收概率则为 $h_1=3/8,h_2=1/2,h_3=5/8$。
\insight 内部漂移指向中间，因此吸收概率通常不是 $k/m$；必须使用真实转移概率构造调和函数。

\problem{R05}{两个重叠模式的竞争}{前缀自动机}
独立信号 H、T 的概率分别为 $p\in(0,1)$、$q=1-p$。从空记录开始，首次出现连续子串 HHT 或 HTH 即止。求 HHT 先出现的概率，并求公平信号时的平均停止时间。
\solution 保留当前后缀中也是目标前缀的最长部分，只需空状态、H、HH、HT 四个非终止状态。转移中，HH 遇 H 仍是 HH，HT 遇 T 回到0。设 HHT 先出现的概率为 $u_s$。有
\[
 u_0=u_H,\quad u_{HH}=1,\quad u_{HT}=qu_0,
 \quad u_H=p+qu_{HT}.
\]
HH 状态下迟早出现 T，因为 $q>0$，所以 $u_{HH}=1$。解得
\[
 \boxed{u_0=\frac1{2-p}}.
\]
公平时记剩余时间为 $e_s$。独立长度3区块出现 HHT 的概率为 $1/8$，保证这些均值有限。首步方程为
\[
 e_0=2+e_H,\quad e_{HH}=2,\quad
 e_{HT}=1+\tfrac12e_0,\quad
 e_H=1+\tfrac12e_{HH}+\tfrac12e_{HT}.
\]
于是 $e_0=9/2+e_0/4$，得 $\boxed{e_0=6}$。
\insight 模式前缀的重叠决定保留哪些信息；两个模式长度相同并不意味着获胜概率相同。

\problem{R06}{马尔可夫信号的连续段计数}{固定窗口与协方差}
平稳二态链取值 H、T，转移概率为 $P(H,T)=a$、$P(T,H)=b$，其中 $0<a,b<1$。观察 $n\ge2$ 个连续信号，相同信号连续出现构成一段。求连续段数 $R_n$ 的均值与方差，并说明 $a=b$ 时的简化。
\solution 令 $I_j=\ind_{\{X_{j+1}\ne X_j\}}$，$m=n-1$，则 $R_n=1+\sum_{j=1}^m I_j$。平稳分布为 $(b/(a+b),a/(a+b))$。置
\[
 r=\frac{2ab}{a+b},\quad \lambda=1-a-b,\quad
 C=\frac{ab(a-b)^2}{(a+b)^2}.
\]
于是 $\E I_j=r$。已知发生切换后，新的状态 H、T 各有一半概率，因为两个方向的平稳流量均为 $ab/(a+b)$。二态链经过 $k-1$ 步后，对下一次切换概率的偏离乘以 $\lambda^{k-1}$，故
\[
 \Prob(I_{j+k}=1\mid I_j=1)
 =r+\lambda^{k-1}\left(\frac{a+b}{2}-r\right),
\]
从而 $\Cov(I_j,I_{j+k})=C\lambda^{k-1}$。求和得
\[
 \boxed{\E R_n=1+mr},
\]
\[
 \boxed{\Var R_n=mr(1-r)
 +2C\frac{m-1-m\lambda+\lambda^m}{(1-\lambda)^2}}.
\]
最后的分式来自 $\sum_{k=1}^{m-1}(m-k)\lambda^{k-1}$，在 $m=1$ 时为0。若 $a=b$，每一步是否切换可由独立的参数 $a$ 伯努利变量实现，故 $R_n-1\sim\operatorname{Bin}(n-1,a)$；这也对应 $C=0$。
\insight 固定窗口中的计数需要协方差，不能仅由单步切换率推断其服从二项分布。

\problem{P15}{移位寄存器的精确混合}{有限记忆}
一个寄存器保留 $r\ge1$ 个二进制位，包括前导零。每步删除最左位、其余左移，并在最右端加入独立比特；新比特为1的概率为 $p\in(0,1)$。以整数表示状态，则 $S_{n+1}=(2S_n+\xi_{n+1})\bmod 2^r$。求平稳分布。若初态全零，求第 $n$ 步到平稳分布的全变差距离。
\solution 经过 $r$ 次更新，全部旧位已经被新比特覆盖，因此状态是最近 $r$ 个独立比特。若 $w(s)$ 为状态 $s$ 的1的个数，则
\[
 \boxed{\pi(s)=p^{w(s)}(1-p)^{r-w(s)}}.
\]
从任意初态在 $r$ 步后都恰有此分布，故它平稳且唯一。对 $n<r$，初态全零时，前 $r-n$ 位固定为0，后 $n$ 位已是独立比特。因此当前分布是 $\pi$ 条件于事件 $A_n=$“前 $r-n$ 位全零”的分布。对概率分布 $\pi$ 和满足 $\pi(A)>0$ 的事件 $A$，条件分布与 $\pi$ 的全变差距离为 $1-\pi(A)$：分别在 $A$ 和补集求绝对概率差即可。于是
\[
 \boxed{d_{\rm TV}(\mathcal L(S_n),\pi)=
 \begin{cases}1-(1-p)^{r-n},&0\le n<r,\\0,&n\ge r.\end{cases}}
\]
\insight 这里遗忘初态是信息被完全覆盖，因此有精确的有限混合时间；平稳分布只有在 $p=1/2$ 时才均匀。


\problem{R07}{单维修员的双机系统}{连续时间链}
两台机器工作时分别以速率 $\lambda>0$ 独立失效。只有一名维修员，每次维修耗时服从参数 $\mu>0$ 的指数分布；只要有坏机便持续维修一台，修好后立即投入工作。所有失效和维修时间独立。至少一台工作时系统可用。求长期可用率，以及从两台全好到首次两台全坏的平均时间。
\solution 以坏机数 $K_t=0,1,2$ 为状态，向上跳率为 $(2-k)\lambda$，向下跳率为 $\mu$。相邻平衡给出
\[
 2\lambda\pi_0=\mu\pi_1,\qquad
 \lambda\pi_1=\mu\pi_2.
\]
令 $D=\mu^2+2\lambda\mu+2\lambda^2$，归一化得
\[
 (\pi_0,\pi_1,\pi_2)=\frac1D(\mu^2,2\lambda\mu,2\lambda^2),
 \qquad \boxed{\text{可用率}=1-\frac{2\lambda^2}{D}}.
\]
有限不可约连续时间链的长期比例由此平稳分布给出。令 $e_0,e_1$ 为到达2前的平均剩余时间。由指数等待和首跳分解，
\[
 e_0=\frac1{2\lambda}+e_1,\qquad
 e_1=\frac1{\lambda+\mu}+\frac\mu{\lambda+\mu}e_0.
\]
整理得
\[
 \boxed{e_0=\frac{3\lambda+\mu}{2\lambda^2}}.
\]
这些均值有限：从任意未吸收状态，在固定长度时间块内发生足够的失效、先于下一次维修完成到达状态2的概率有统一正下界。
\insight 两台坏机共享维修员时，总修复速率仍是 $\mu$，不能将两台机器当成独立的双态系统。

\problem{R08}{加权三角形的首达时间}{首步方程}
三角形顶点 A、B、C 的边权分别为 $w_{AB}=a,w_{BC}=b,w_{CA}=c$，均为正。每步按相邻边权占比选择下一顶点。求从 A 到 B、从 B 到 A 的首达时间均值及往返均值。固定 $a=b=1$ 时，增大 $c$ 是否加快从 A 到 B？
\solution 令 $u_A,u_C$ 为命中 B 的均值。首步分解为
\[
 u_A=1+\frac c{a+c}u_C,\qquad
 u_C=1+\frac c{b+c}u_A.
\]
有限不可约链保证均值有限。置 $D=ab+ac+bc$，消去 $u_C$ 得
\[
 \boxed{\E_A T_B=\frac{(a+2c)(b+c)}D}.
\]
互换 A、B 的角色，即互换 $b,c$。若 $C$ 为先到 B 再回 A 的总步数，由强马尔可夫性，得到
\[
 \boxed{\E_B T_A=\frac{(a+2b)(b+c)}D},\qquad
 \boxed{\E_A C
 =\frac{2(a+b+c)(b+c)}D}.
\]
若 $a=b=1$，第一式化为 $\boxed{1+c}$，所以增大 A--C 边权反而使到 B 更慢：游走越来越多地在 A、C 间往返。
\insight 加大一条边的权重也会降低其他出口被选择的概率；加权随机游走的步数不能只按图距离判断。

\problem{R09}{反射线段的覆盖时间}{边界分解}
随机游走位于 $\{0,1,\ldots,N\}$，其中 $N\ge2$。内部每步等概率增减1；在0下一步必到1，在 $N$ 下一步必到 $N-1$。从 $1\le i<N$ 出发，首次访问过两个端点时停止，时间为 $C$。求最后访问的是 $N$ 的概率与 $\E_iC$。
\solution 先到任一端点的时间记为 $\sigma$。到达端点前是普通区间吸收游走，其调和方程及时间方程给出
\[
 \Prob_i(X_\sigma=0)=\frac{N-i}{N},\qquad
 \E_i\sigma=i(N-i).
\]
例如时间函数 $v_i=i(N-i)$ 满足 $v_i=1+(v_{i-1}+v_{i+1})/2$ 和零边界。最后访问 $N$ 等价于先访问0，所以第一项概率即为所求。

令 $u_j$ 为从 $j$ 到 $N$ 的均值，保留0的反射规则。则 $u_N=0$、$u_0=1+u_1$，内部满足同一首步时间方程。二次多项式 $u_j=N^2-j^2$ 满足全部条件，故从0横穿到 $N$ 的均值为 $N^2$。从 $N$ 到0同理。利用强马尔可夫性，
\[
 \boxed{\Prob_i(\text{末访 }N)=\frac{N-i}{N}},\qquad
 \boxed{\E_i C=i(N-i)+N^2}.
\]
有限状态链具有统一正概率在有限步横穿区间，保证所用停止时间可积；访问两个端点也必访问所有中间点。

\problem{R10}{成对抽样的覆盖}{容斥与尾和}
有 $n\ge3$ 种标签，每轮独立、均匀抽取一个由两种不同标签组成的无序对。首次见过全部标签时停止，轮数为 $T$。给出 $T$ 的尾概率、均值及二阶矩的有限和表达式，并计算 $n=4$ 时的均值与方差。
\solution 对一个含 $k$ 种标签的指定集合，一轮完全避开它的概率为
\[
 q_k=\frac{\binom{n-k}{2}}{\binom n2},
\]
其中上标小于2时组合数取0。经过 $t$ 轮仍未覆盖，等价于至少一种标签未出现。容斥给出
\[
 \Prob(T>t)=\sum_{k=1}^n(-1)^{k+1}\binom nk q_k^t,
 \qquad t=0,1,\ldots,
\]
在 $t=0$ 约定 $0^0=1$。所有 $q_k<1$，故可对有限项逐项求几何级数。利用整数尾和恒等式，
\[
 \boxed{\E T=\sum_{k=1}^n(-1)^{k+1}\binom nk\frac1{1-q_k}},
\]
\[
 \boxed{\E T^2=\sum_{k=1}^n(-1)^{k+1}\binom nk
 \frac{1+q_k}{(1-q_k)^2}}.
\]
当 $n=4$，$(q_1,q_2,q_3,q_4)=(1/2,1/6,0,0)$，代入得到
\[
 \boxed{\E T=\frac{19}{5}},\quad
 \E T^2=\frac{423}{25},\quad
 \boxed{\Var T=\frac{62}{25}}.
\]
\insight 一轮可能发现两个新标签，不能把各发现阶段当作依次独立的几何等待；“哪些标签尚未出现”的事件更适合容斥。

\par\Needspace{8\baselineskip}
\subsection*{小节练习}
\exercise{OR01}{连续正面的等待时间}
独立重复抛掷一枚公平硬币。设 $n\ge1$ 为整数，首次连续出现 $n$ 次正面时停止。求总抛掷次数 $T_n$ 的期望。

\exercise{OR07}{雨伞的长期短缺频率}
某人共有 $n\ge1$ 把雨伞，初始家中有 $k$ 把、办公室有 $n-k$ 把，其中 $0\le k\le n$。他每天早晨从家去办公室，晚上返回。每次出发时独立地以概率 $p\in[0,1]$ 下雨；下雨且出发地有伞时带走一把，不下雨时不带伞。下雨而出发地没有伞时会淋雨。求长期每次出行的淋雨频率。


\section{随机游走}
\noindent\textbf{要点与方法.}\enspace
随机游走须先写清吸收、反射和重启规则；方向有记忆或边界随历史移动时，应把方向或回撤差距纳入状态。命中概率满足调和方程，平均时间满足带常数项的首步方程。对成功路径作条件化时，要按后续成功概率重新加权，不能只删去失败终点。

鞅提供另一类边值工具，但无界停时不能直接代入期望恒等式。可先停在 $\tau\wedge n$，再凭一致可积性或可积支配取极限；仅有几乎必然停止并不够。跳跃可能越过边界，终值也须按真实越界位置处理。
\par\smallskip



\problem{R11}{带重启的清空时间}{重置与最优化}
整数过程从1出发，到0即止。每个未停止时刻，以概率 $\rho\in(0,1)$ 将位置重置到1；否则等概率向左或向右走一步。重置也计一步，所有选择独立。求停止时间均值，并选择使它最小的 $\rho$。
\solution 从任意正位置，连续两步“重置、向左”以概率 $\rho(1-\rho)/2$ 结束，故均值有统一有限上界。记从 $i$ 出发的均值为 $e_i$，则
\[
 e_i=1+\rho e_1+\frac{1-\rho}{2}(e_{i-1}+e_{i+1}),
 \qquad e_0=0.
\]
特征根互为倒数，取小于1的根
\[
 r=\frac{1-\sqrt{1-(1-\rho)^2}}{1-\rho}.
\]
常数特解为 $(1+\rho e_1)/\rho$；有界性排除大根，零边界给出
\[
 e_i=\frac{1+\rho e_1}{\rho}(1-r^i),\qquad
 \boxed{e_1=\frac{1-r}{\rho r}}.
\]
由根方程 $1-\rho=2r/(1+r^2)$，改以 $r\in(0,1)$ 为变量，
\[
 \rho=\frac{(1-r)^2}{1+r^2},\qquad
 e_1=\frac{1+r^2}{r(1-r)}.
\]
导数的符号由 $r^2+2r-1$ 决定，两端的均值都发散，故唯一最优点为 $r=\sqrt2-1$。因此
\[
 \boxed{\rho_*=1-\frac1{\sqrt2}},\qquad
 \boxed{\min\E T=2+2\sqrt2}.
\]
\insight 重启既能拉回远离目标的轨迹，也会浪费已接近目标的时间；最优重启率来自这两种效应的权衡。

\problem{R12}{成功吸收路径的占用次数}{条件链与Green函数}
公平整数游走从 $i\in\{1,\ldots,N-1\}$ 出发，到0或 $N$ 即止，时间为 $\tau$。已知最终到达 $N$，求每个内部点 $j$ 在停止前的期望访问次数（包括初始时刻），并由此求条件平均时间。
\solution 先求无条件占用函数
\[
 G(i,j)=\E_i\sum_{t=0}^{\tau-1}\ind_{\{X_t=j\}}.
\]
固定 $j$，它在 $i\ne j$ 处满足线性调和方程，在 $j$ 处满足 $G(j,j)=1+[G(j-1,j)+G(j+1,j)]/2$，两端为0。左右分别取一次函数并匹配 $j$ 处的斜率跳跃，得到
\[
 G(i,j)=\frac{2\min(i,j)[N-\max(i,j)]}{N}.
\]
成功概率为 $h_i=i/N$。在一次访问 $j$ 后，未来成功概率为 $h_j$；对非负访问计数逐项用条件期望，得到成功事件加权的占用均值 $h_jG(i,j)$。除以 $h_i$，
\[
 \boxed{G^*(i,j)=\frac{2j\min(i,j)[N-\max(i,j)]}{Ni}}.
\]
按 $j\le i$ 与 $j>i$ 分组求和，并使用平方和公式，得到
\[
 \boxed{\E_i(\tau\mid X_\tau=N)=\sum_{j=1}^{N-1}G^*(i,j)
 =\frac{N^2-i^2}{3}}.
\]
也可检验未除以成功概率的函数 $v_i=i(N^2-i^2)/(3N)$ 满足 $v_i=h_i+(v_{i-1}+v_{i+1})/2$ 和零边界。有限区间的几何退出尾保证上述期望有限。
\insight 成功条件将对点 $j$ 的访问按 $h_j/h_i$ 重加权，而不是把不成功的终点简单删掉后仍保留原转移概率。

\problem{R13}{带方向记忆的区间游走}{扩充状态}
游走在 $\{0,\ldots,N\}$ 上，到端点即止，$N\ge2$。内部每步以概率 $\alpha\in(0,1)$ 保持上一步方向，以概率 $1-\alpha$ 反向，然后按所选方向移动一格。初始位置为 $1\le i<N$，并已知上一步方向。求前一方向分别为右、左时，先到 $N$ 的概率 $h_i^+,h_i^-$。
\solution 位置本身不构成马尔可夫状态，须同时记方向。内部方程为
\[
 h_i^+=\alpha h_{i+1}^++(1-\alpha)h_{i-1}^-,\qquad
 h_i^-=(1-\alpha)h_{i+1}^++\alpha h_{i-1}^-.
\]
实际可能进入的吸收状态只有 $h_N^+=1,h_0^-=0$。尝试 $h_i^-=Ai$、$h_i^+=Ai+B$，代入得 $(1-\alpha)B=(2\alpha-1)A$；再用 $AN+B=1$，令
\[
 D=N(1-\alpha)+2\alpha-1=(N-2)(1-\alpha)+1>0,
\]
便有
\[
 \boxed{h_i^- =\frac{i(1-\alpha)}D,\qquad
 h_i^+=\frac{i(1-\alpha)+2\alpha-1}D}.
\]
这些候选值满足全部方程。连续选择向右有统一正概率在至多 $N$ 步吸收，故有限状态边值问题的概率解唯一。$\alpha=1/2$ 时两式都化为 $i/N$；$\alpha>1/2$ 时，向右的初始方向提高成功概率。
\insight 只对左右步数做平均会丢失方向持续性；初始方向是题设信息的一部分。

\problem{R14}{给定终点的三角形游走}{马尔可夫桥}
三角形 A、B、C 上每步等概率走向另外两个顶点，从 A 开始。给定第 $n\ge2$ 步位于 A，对整数 $0\le k\le n$，求第 $k$ 步位于 B 的概率，并求从时刻0到 $n$ 访问 A 的总次数均值。计算 $n=4,k=2$。
\solution 置 $r=-1/2$。由 $a_{t+1}=(1-a_t)/2,a_0=1$，返回概率及指定异点概率为
\[
 a_t=\frac{1+2r^t}{3},\qquad b_t=\frac{1-r^t}{3}.
\]
使用时刻 $k$ 的马尔可夫性及转移对称性，
\[
 \boxed{\Prob(X_k=B\mid X_n=A)=\frac{b_kb_{n-k}}{a_n}}.
\]
若 $L=\sum_{k=0}^n\ind_{\{X_k=A\}}$，则
\[
 \E(L\mid X_n=A)=\frac1{a_n}\sum_{k=0}^n a_ka_{n-k}.
\]
将乘积展开为 $[1+2r^k+2r^{n-k}+4r^n]/9$ 并求几何级数，得到
\[
 \boxed{\E(L\mid X_n=A)=
 \frac{(n+1)(1+4r^n)+4(1-r^{n+1})/(1-r)}{3(1+2r^n)}}.
\]
当 $n=4,k=2$，$a_4=3/8,b_2=1/4$，所求概率为 $\boxed{1/6}$，访问 A 的均值为 $\boxed{8/3}$。条件链通常不齐次：在时刻 $k$ 从 $u$ 到 $v$ 的概率应乘以 $P^{n-k-1}(v,A)/P^{n-k}(u,A)$，仅在分母为正的条件可达状态使用。

\problem{R15}{轮流累加的越界游戏}{更新次数的奇偶性}
两人轮流产生独立的 $U[0,1]$ 增量，全部增量累加到同一个初始为0的总量中。首次使总量超过 $t\in(0,1]$ 者获胜。求先手胜率、使游戏公平的 $t$，以及总轮数 $N$ 的均值与方差。
\solution 对 $n\ge1$，未越界等价于前 $n$ 个增量的和不超过 $t$。由于 $t\le1$，该区域是边长约束未起作用的单纯形，体积为
\[
 \Prob(N>n)=\frac{t^n}{n!};\qquad \Prob(N>0)=1.
\]
尾概率衰减保证所有矩有限。概率母函数可由尾和恒等式写为
\[
 G(z)=\E z^N=1+(z-1)\sum_{n=0}^\infty z^n\Prob(N>n)
 =1+(z-1)e^{zt}.
\]
先手获胜当且仅当 $N$ 为奇数，因此
\[
 \boxed{\Prob(\text{先手胜})=\frac{1-G(-1)}2=e^{-t}},\qquad
 \boxed{t_{\rm fair}=\log2}.
\]
求导得 $G'(1)=e^t,G''(1)=2te^t$，于是
\[
 \boxed{\E N=e^t},\qquad
 \boxed{\Var N=(2t+1)e^t-e^{2t}}.
\]
\insight 单调累计把胜负化为停止次数的奇偶性；仅知道平均轮数不能推断先手优势。

\problem{R16}{相对历史最高点的回撤}{移动边界}
公平整数游走 $S_n$ 从0出发，令 $M_n=\max_{0\le j\le n}S_j$。给定整数 $d\ge1$，首次出现 $M_n-S_n=d$ 时停止，时间为 $\tau$。求 $\E\tau,\Var\tau$，并给出停止前最高点 $M_\tau$ 的分布。
\solution 差距 $Y_n=M_n-S_n$ 是有限状态链：在0以各 $1/2$ 的概率留在0或到1，在 $1,\ldots,d-1$ 等概率增减1，在 $d$ 吸收。连续 $d$ 次向下使它结束，故所有停止矩有限。令 $m_i=\E_i\tau$，则 $m_0=1+(m_0+m_1)/2$，内部满足 $m_i=1+(m_{i-1}+m_{i+1})/2$，$m_d=0$。二次解为
\[
 m_i=D-i(i+1),\qquad D=d(d+1),\qquad \boxed{\E\tau=D}.
\]
令 $v_i$ 为二阶矩。展开 $(1+\tau')^2$ 给出
\[
 v_0-v_1=4D-2,\qquad
 v_{i-1}-2v_i+v_{i+1}=2-4m_i,\qquad v_d=0.
\]
对一阶差分逐次求和，
\[
 v_0=\sum_{j=0}^{d-1}(d-j)(4m_j-2)=\frac{5D^2-D}{3},
 \qquad \boxed{\Var\tau=\frac{D(2D-1)}3}.
\]
为求最高点，在每次新高 $h$ 处重新开始：先到 $h+1$ 而不是 $h-d$ 的概率为 $d/(d+1)$。用强马尔可夫性重复此竞争，
\[
 \boxed{\Prob(M_\tau\ge k)=\left(\frac d{d+1}\right)^k,
 \quad k=0,1,\ldots}.
\]
\insight 绝对位置与最高点组成二维状态，而差距将移动边界变成带自环的反射边界。

\problem{R17}{固定比例投注的财富极限}{非负鞅与越界}
财富从 $x>0$ 开始，每轮用财富的固定比例 $f\in(0,1)$ 公平投注，赢则乘以 $1+f$，输则乘以 $1-f$，各轮独立。记财富为 $Z_n$。证明 $\E Z_n=x$，但 $Z_n\to0$ 几乎必然。对目标 $A>x$，给出曾达到或超过 $A$ 的概率上下界，并解释为何不能把越界财富直接设为 $A$。
\solution 因为条件平均乘数为1，$Z_n$ 是非负鞅，且每个有限时刻 $\E Z_n=x$。另一方面，由强大数律，
\[
 \frac1n\log\frac{Z_n}{x}\longrightarrow
 \tfrac12\log(1+f)+\tfrac12\log(1-f)
 =\tfrac12\log(1-f^2)<0.
\]
故财富几乎必然趋0。其均值不随之趋0，说明该鞅族不一致可积。

令 $\tau=\inf\{n:Z_n\ge A\}$。对有界停时 $n\wedge\tau$，$\E Z_{n\wedge\tau}=x$。停止族被常数 $(1+f)A$ 控制：未命中时低于 $A$，命中时的上一步仍低于 $A$。其极限在命中事件上为 $Z_\tau$，在未命中事件上为0。支配收敛得
\[
 x=\E[Z_\tau\ind_{\{\tau<\infty\}}].
\]
利用 $A\le Z_\tau<(1+f)A$，
\[
 \boxed{\frac{x}{(1+f)A}\le\Prob(\tau<\infty)\le\frac xA}.
\]
\insight 公平的算术收益不保证典型路径增长。离散乘法跳跃会越过目标，终端财富不是固定边界值。

\problem{R18}{首次碰撞的累计成本}{停止抽样与选择偏差}
每次独立抽取类型 $j\in\{1,\ldots,m\}$，概率 $p_j>0$，总和为1，费用为 $c_j\ge0$。首次抽到此前见过的类型时停止，次数为 $N$，含最后一次。求期望总费用，并求最后一次类型的分布。计算两类概率 $2/3,1/3$、费用1和4的情形。
\solution 记 $e_k(p)$ 为 $p_1,\ldots,p_m$ 的第 $k$ 个基本对称多项式，$e_0=1$。前 $k$ 次互不相同的概率为 $k!e_k(p)$，所以
\[
 \boxed{\E N=\sum_{k=0}^m k!e_k(p)}.
\]
这里 $2\le N\le m+1$。设第 $t$ 次费用为 $C_t$，$\bar c=\sum_jp_jc_j$。事件 $N\ge t$ 只依赖前 $t-1$ 次，与 $C_t$ 独立；有限求和得到
\[
 \boxed{\E\sum_{t=1}^N C_t=\bar c\,\E N}.
\]
若最后重复类型为 $j$ 且 $N=k+1$，前 $k$ 个互异类型必须包含 $j$。将这些类型排序并再抽一次 $j$，其概率为 $k!p_j^2e_{k-1}(p_{-j})$，其中 $p_{-j}$ 表示删去第 $j$ 项。因此
\[
 \boxed{\Prob(J_{\rm last}=j)=p_j^2\sum_{k=1}^m k!e_{k-1}(p_{-j})}.
\]
两类情形置 $p=2/3,q=1/3$，有 $\E N=2+2pq=22/9$，$\bar c=2$，故总费用均值为 $\boxed{44/9}$。最后是第一类的概率为 $p^2(1+2q)=20/27$，最后一次费用的均值为 $\boxed{16/9}$，并非单次无条件均值2。
\insight 累计费用可以用过去可测的继续事件分解；最后一次的类型却受到停止规则的选择，不能套同一个独立性结论。

\par\Needspace{8\baselineskip}
\subsection*{小节练习}
\exercise{OR11}{非对称游走的命中概率}
随机游走从正整数 $n$ 出发，每步独立地以概率 $p\in(0,1)$ 向左走1，以概率 $1-p$ 向右走2。求此后至少一次到达0的概率 $u_n$。

\exercise{OR16}{双边吸收时间的二阶矩}
设 $S_0=0$，$S_n$ 为每步独立、等概率增减1的随机游走。给定正整数 $a,b$，令
\[
 \tau=\inf\{n\ge0:S_n=-a\text{ 或 }S_n=b\}.
\]
求 $\Prob(S_\tau=b)$、$\E\tau$ 及 $\E\tau^2$。


\section{最优停止}
\noindent\textbf{要点与方法.}\enspace
最优停止比较的是当前信息下“立即停止”与“继续后采用最优策略”的价值。状态应包括剩余机会、可召回报价或未知参数的后验；已支付的费用不应再次扣除。先逐点比较可行动作，再对观察结果取期望，等号处可能有多种最优决策。

有限期限可用反向归纳证明策略最优。无限期限则须验证价值有限、递推极限可靠且候选策略达到该值；折扣可给出统一的截断误差界。涉及复制与超复制时，还须逐节点检查自融资、借款和卖空约束，不能只比较终端支付。
\par\smallskip



\problem{R19}{融资约束下的数字支付}{复制与超复制}
股票初价100，两期内每期价格乘以2或 $1/2$，每条路径都有正概率；现金账户利率为0，可交易任意实数份股票且无摩擦。期末支付为 $\ind_{\{S_2\ge100\}}$。先允许借款、卖空，求精确复制的初始资金与逐期持仓。再要求每次调仓后股票份数和现金均非负，求能覆盖所有路径支付的最小初始资金。
\solution 一期后若股价为200，两种期末价格的支付均为1，持现金1即可。若股价为50，后继价格为100、25，复制组合 $(\Delta,B)$ 满足
\[
 100\Delta+B=1,\qquad25\Delta+B=0,
\]
故 $\Delta=1/75,B=-1/3$，此节点价值为 $1/3$。初始组合满足 $200\Delta_0+B_0=1$、$50\Delta_0+B_0=1/3$，从而
\[
 \boxed{\Delta_0=\frac1{225},\quad B_0=\frac19,\quad V_0=\frac59}.
\]
到达任一节点时旧组合的价值恰好等于新组合成本，因此策略自融资。

现在施加非负持仓约束。低节点中，为覆盖价格100时的支付，需要 $100\Delta+B\ge1$，故成本 $50\Delta+B\ge1/2$；取 $\Delta=1/100,B=0$ 达到。高节点的两个支付都为1，约束 $100\Delta+B\ge1$ 给出成本 $200\Delta+B\ge1$，现金1达到。因此初始须覆盖节点资金1、$1/2$。对任意可行组合，
\[
 100\Delta_0+B_0
 =\tfrac13(200\Delta_0+B_0)+\tfrac23(50\Delta_0+B_0)
 \ge\frac23.
\]
取 $\Delta_0=1/300,B_0=1/3$，两个节点资金恰为1、$1/2$，可按上述策略继续，故
\[
 \boxed{V_0^{\rm super}=\frac23}.
\]
\insight 精确复制中的低节点需要借款。禁止借款后，应逐节点求超复制成本，不能仍把无约束价格当成可实施资金。

\problem{R20}{可召回报价的停止决策}{有限机会}
先免费观察 $X\sim U[0,1]$。可接受 $X$，或支付 $c>0$ 再观察独立的 $Y\sim U[0,1]$，随后取得 $\max(X,Y)$。最多再观察一次。求最优策略与免费开始的最优净价值。取 $c=1/8$，并与支付后必须接受 $Y$、不能召回 $X$ 的规则比较。
\solution 给定 $X=x$，追加观察相对停止的期望增益为
\[
 \E(Y-x)^+-c=\int_x^1(y-x)\dd y-c=\frac{(1-x)^2}{2}-c.
\]
所以当 $0<c<1/2$ 时，仅在 $x<1-\sqrt{2c}$ 时继续，等号时两者同优；$c\ge1/2$ 时始终停止。最优价值为
\[
 V=\frac12+\int_0^{1-\sqrt{2c}}
 \left(\frac{(1-x)^2}{2}-c\right)\dd x
 =\frac23-c+\frac{(2c)^{3/2}}3
\]
（此式用于 $c<1/2$；否则 $V=1/2$）。因此 $c=1/8$ 时门槛为 $1/2$，$\boxed{V=7/12}$。

不能召回时，继续价值固定为 $\E Y-c=3/8$。于是只在 $x<3/8$ 时继续，净价值为
\[
 \boxed{\E\max(X,3/8)=\frac12+\frac12\left(\frac38\right)^2
 =\frac{73}{128}}.
\]
前者较高；召回权保留了旧报价的下限，改变了观察新信息的价值。

\problem{R21}{学习未知胜率后的停止}{贝叶斯决策}
先从 $U[0,1]$ 抽取一个未知参数 $P$，随后给定 $P$ 独立抛币，H 概率为 $P$。每次 H 收益 $+1$、T 收益 $-1$，可在任意时刻停止，最多抛三次，也可直接停止。已取得的收益保留。求最大期望总收益和完整最优规则，并与已知公平硬币比较。
\solution 已见 $h$ 个 H、$t$ 个 T 时，参数后验密度正比于 $p^h(1-p)^t$，下一次 H 的概率为 $q_{h,t}=(h+1)/(h+t+2)$。令 $V_k(h,t)$ 为还可抛 $k$ 次的最优未来增益。则 $V_0=0$，且
\[
 \begin{aligned}
 C_k(h,t)&=q_{h,t}[1+V_{k-1}(h+1,t)]\\
 &\quad +(1-q_{h,t})[-1+V_{k-1}(h,t+1)],\\
 V_k(h,t)&=\max\{0,C_k(h,t)\}.
 \end{aligned}
\]
停止取得未来增益0；继续项已经对下一结果及后续最优决策取期望，因而反向归纳证明该递推最优。最后一次机会时，
\[
 V_1(2,0)=\tfrac12,\qquad V_1(1,1)=V_1(0,2)=0.
\]
进一步得到
\[
 V_2(1,0)=\tfrac23,\quad V_2(0,1)=0,\qquad
 \boxed{V_3(0,0)=\tfrac13}.
\]
因此先抛一次；若为 T 就停止，若为 H 则再抛。之后 HH 时继续，HT 时最后一次的条件均值为0，停止或继续都最优。若硬币事先已知公平，任意至多三次的停止策略，其有界停时收益均值均为0。
\insight 每次结果的边际概率虽为 $1/2$，观测后预测概率却改变；未知的共享参数使多次结果并不独立。

\problem{R22}{折扣报价的无限期停止}{压缩映射}
报价 $X_1,X_2,\ldots$ 独立服从 $U[0,1]$，观察后可接受或永久放弃该报价。第 $N$ 次接受的现值为 $\delta^{N-1}X_N$，其中 $0<\delta<1$；永不接受的支付为0。求最优停止规则、初始价值及采用严格门槛时的平均观察次数。
\solution 最多可观察 $k$ 次的价值记为 $V_k$，$V_0=0$。反向归纳给出
\[
 V_{k+1}=\E\max(X,\delta V_k).
\]
映射的 Lipschitz 常数至多为 $\delta<1$，所以迭代收敛到唯一不动点 $V$。任意无限期策略截在第 $k$ 次，支付改变的绝对值至多 $\delta^{k-1}$，因此不动点就是无限期最优值。

令 $t=\delta V\in(0,1)$，直接积分得 $V=(1+t^2)/2$，即 $\delta t^2-2t+\delta=0$。选择 $(0,1)$ 内的根，
\[
 \boxed{t=\frac{1-\sqrt{1-\delta^2}}\delta},\qquad
 \boxed{V=\frac1{1+\sqrt{1-\delta^2}}}.
\]
首次 $X_n>t$ 时接受，观察次数为参数 $1-t$ 的几何变量，故 $\boxed{\E N=1/(1-t)}$。该策略的实际现值为
\[
 \sum_{n=1}^\infty(\delta t)^{n-1}\E[X\ind_{\{X>t\}}]
 =\frac{1-t^2}{2(1-\delta t)}=V,
\]
验证其达到上界。
\insight 折扣使远期支付误差有确定的几何上界，因而无限期问题可以由有限期递推严格逼近。

\problem{R23}{两枚骰子的重掷选择}{二维观察与控制}
同时掷两枚独立公平六面骰，观察点数后可直接取得乘积，或支付 $c\ge0$ 选择其中一枚重掷一次。重掷后必须采用新点数，取得最终乘积减费用。求完整最优规则，并计算 $c=1$ 的最大期望净收益。
\solution 将已见点数排序为 $x\le y$。重掷较小骰的净期望为 $(7/2)y-c$，重掷较大骰为 $(7/2)x-c$，前者不低于后者。因此最优规则为
\[
 \boxed{\text{重掷较小骰}\quad\Longleftrightarrow\quad
 y(7/2-x)>c},
\]
等号可任选，条件不满足时停止。无控制时的乘积均值为 $(7/2)^2=49/4$，故最优价值亦可写为
\[
 V(c)=\frac{49}{4}+\frac1{36}\sum_{i=1}^6\sum_{j=1}^6
 \left(\tfrac72\max(i,j)-ij-c\right)^+.
\]
当 $c=1$，继续恰好发生在较小点数不超过3时。按较小点数为1、2、3分组，36个等概率有序结果中的增益和分别为 $183/2,48,19/2$；相等点数只计一次，其余计两次。因此
\[
 \boxed{V(1)=\frac{49}{4}+\frac{149}{36}=\frac{295}{18}}.
\]
\insight 应先比较两个可控动作，再与停止比较；“只要出现小点数就重掷”的规则还必须经过费用检验。

\problem{R24}{带早停的测试顺序}{相邻交换}
独立测试 $i=1,\ldots,n$ 的通过概率为 $p_i\in(0,1)$，执行成本为 $c_i>0$。为判断是否全部通过，按某个可自适应选择的顺序测试，遇到第一个失败就停止；全部通过则做完。求最小期望成本的顺序。若目标改为判断是否至少一项通过，规则应如何变化？计算 $p=(9/10,1/2,1/10)$、$c=(1,2,5)$。
\solution 全通过判定中，只要尚未停止，历史结果必定全通过。因此任一确定性自适应策略等价于一个固定排列；随机化只是排列的混合。排列 $\pi$ 的期望成本为
\[
 \sum_{k=1}^n c_{\pi_k}\prod_{j<k}p_{\pi_j}.
\]
相邻测试 $i,j$ 的局部成本分别为 $c_i+p_ic_j$ 和 $c_j+p_jc_i$。两者之后继续的概率均为 $p_ip_j$，故后续成本不变。将较便宜者放前等价于
\[
 c_i(1-p_j)\le c_j(1-p_i)
 \quad\Longleftrightarrow\quad
 \frac{c_i}{1-p_i}\le\frac{c_j}{1-p_j}.
\]
逐次交换逆序对，得到按 $c_i/(1-p_i)$ 升序的全局最优排列。

至少一项通过的判定在首次通过时停止，未停止历史全为失败。同一交换证明将继续概率换为 $1-p_i$，因此按 $c_i/p_i$ 升序。给定数据下，两种最优顺序和成本为
\[
 \boxed{\text{全部通过： }(2,3,1),\quad C=\frac{91}{20}}.
\]
\[
 \boxed{\text{至少一项： }(1,2,3),\quad C=\frac{29}{20}}.
\]
\insight 排序指标包含停止目标；同一组测试，在不同逻辑目标下可以有不同的最优顺序。

\par\Needspace{8\baselineskip}
\subsection*{小节练习}
\exercise{OR20}{红黑牌的最优停止}
一副牌含26张红牌和26张黑牌，均匀洗牌后逐张抽取，不放回。初始得分为0，每抽到一张红牌加1分，黑牌减1分。可以在尚未抽牌时或任意一次抽牌后停止，得到当前累计分数。求使期望得分最大的停止规则及最大期望得分。

\exercise{OR22}{付费重掷骰子}
一枚公平六面骰子的点数为 $1,2,\ldots,6$。第一次投掷免费；看到点数后，可以接受并结束，也可以支付1元重掷，之后每次重掷都支付1元。所有投掷相互独立，不能保留先前点数。策略须几乎必然在有限次投掷后结束。最终所得为最后一次点数减去累计重掷费用。求最优策略及最大期望所得。


\section{布朗运动}
\noindent\textbf{要点与方法.}\enspace
布朗运动的独立增量与高斯结构可把路径计算化为协方差积分。确定的平方可积核对应高斯伊藤积分；指数化或在独立随机时刻取值后，一般不再正态。联合高斯条件可用正交投影处理，最大值可用反射原理；对终点等零概率条件，须通过密度或连续条件版本解释。

独立指数失效可转成贴现因子，再用生成元边值方程和伊藤公式验证。交换积分、期望或参数导数时，应说明非负性、可积性或支配条件。泊松分流要求独立标记；更新过程在长时间区间内均匀选点观察，其极限分布须按间隔长度加权。
\par\smallskip



\problem{R25}{指数布朗积分的预测}{非高斯投影}
设 $W$ 为从0出发的标准布朗运动，$T>0$，令
\[
 I=\int_0^T e^{W_t-t/2}\dd t.
\]
求 $I$ 的均值、方差，及用 $\alpha+\beta W_T$ 预测 $I$ 的最优系数和均方误差。再求给定 $W_T=b$ 的条件均值，判断它是否等于最优仿射预测。
\solution 记 $M_t=e^{W_t-t/2}$。正态矩母函数给出 $\E M_t=1$，且当 $s\le t$，$\E M_sM_t=e^s$。由非负被积函数的 Tonelli 定理，
\[
 \E I=T,\qquad
 \E I^2=2\int_0^T\int_0^t e^s\dd s\dd t
 =2(e^T-1-T).
\]
结果有限，故后续协方差交换也由平方可积性保证。因此
\[
 \boxed{\Var I=2(e^T-1-T)-T^2}.
\]
独立增量及正态矩母函数求导给出 $\E(M_tW_T)=\E(M_tW_t)=t$，所以 $\Cov(I,W_T)=T^2/2$。二次损失的正规方程得到
\[
 \boxed{\alpha=T,\qquad\beta=T/2}.
\]
\[
 \boxed{\operatorname{MSE}_{\rm affine}=2(e^T-1-T)-T^2-\frac{T^3}{4}}.
\]
给定 $W_T=b$ 时，$W_t$ 为均值 $tb/T$、方差 $t-t^2/T$ 的正态变量；这来自 $W_t-(t/T)W_T$ 与 $W_T$ 的高斯正交分解。再次使用正态指数矩及条件 Tonelli，
\[
 \boxed{\E(I\mid W_T=b)=\int_0^T
 \exp\left(\frac{tb}{T}-\frac{t^2}{2T}\right)\dd t}.
\]
该函数对 $b$ 的二阶导数严格为正，故不是仿射函数。允许非线性预测可严格降低均方误差。
\insight 指数化破坏了联合高斯结构；“协方差给出最优线性系数”不等于“条件均值一定线性”。

\problem{R26}{给定终点的路径最大值}{反射原理}
设 $W$ 为标准布朗运动，$M_T=\max_{0\le t\le T}W_t$，$T>0$。给定 $W_T=b$，求 $M_T$ 的条件分布。特别地，求 $b=0$ 时最大值的均值。
\solution 若 $a>\max(0,b)$，把首次到 $a$ 以后的路径关于水平线 $a$ 反射。此映射将“曾越过 $a$ 且终点在 $b$ 附近”的路径对应到终点在 $2a-b$ 附近的路径。因此以 $\varphi_T$ 表示 $N(0,T)$ 密度，条件尾概率为
\[
 \boxed{\Prob(M_T\ge a\mid W_T=b)
 =\frac{\varphi_T(2a-b)}{\varphi_T(b)}
 =\exp\left(-\frac{2a(a-b)}T\right)}.
\]
这是密度之比给出的连续版本的正则条件分布，不是对零概率事件直接取比例。若 $a\le\max(0,b)$，尾概率为1。

当 $b=0$，最大值非负，其分布函数为 $1-e^{-2a^2/T}$，$a\ge0$。由非负随机变量的尾积分，
\[
 \boxed{\E(M_T\mid W_T=0)=\int_0^\infty e^{-2a^2/T}\dd a
 =\sqrt{\frac{\pi T}{8}}}.
\]
\insight 终点约束不等于整条路径约束；即使终点回到0，路径仍可能越过正障碍，反射法计算的正是这部分概率。

\problem{R27}{终点与面积双条件的布朗路径}{高斯条件化}
设 $W$ 为标准布朗运动，$T>0$，$A=\int_0^T W_t\dd t$。给定 $W_T=b$ 且 $A=a$，求整个条件过程的均值函数和协方差函数，并说明为何两处路径值可以出现负的条件协方差。
\solution 先去掉终点方向，令 $B_t=W_t-(t/T)W_T$。由联合高斯性，$B$ 与 $W_T$ 独立，且
\[
 K(s,t)=\Cov(B_s,B_t)=\min(s,t)-\frac{st}{T}.
\]
再令 $J=\int_0^T B_t\dd t=A-TW_T/2$。有界连续协方差保证积分平方可积，逐次积分得
\[
 \Cov(B_t,J)=\frac{t(T-t)}2,\qquad \Var J=\frac{T^3}{12}.
\]
将 $B_t$ 投影到 $J$ 上，残差与 $J,W_T$ 独立；该结论对任意有限个时刻同时成立。因此条件过程仍为高斯过程，均值为
\[
 \boxed{m(t)=\frac{tb}{T}
 +\frac{6t(T-t)}{T^3}\left(a-\frac{Tb}{2}\right)},
\]
协方差为
\[
 \boxed{K_*(s,t)=\min(s,t)-\frac{st}{T}
 -\frac{3s(T-s)t(T-t)}{T^3}}.
\]
两个观测的协方差矩阵非退化，因为 $\Var J>0$，所以这些公式给出所有实数 $a,b$ 的连续条件版本。取 $s=T/4,t=3T/4$，有 $K_*(s,t)=-11T/256<0$。固定总面积后，一处偏高可能由另一处偏低补偿；这与原布朗运动的正协方差不同。

\problem{R28}{失效时钟下的区间首达}{贴现边值问题}
过程 $X_t=x+\sigma W_t$ 从 $0<x<L$ 出发，$\sigma>0$，首次到0或 $L$ 时刻为 $\tau$。另有与整条路径独立的指数失效时钟 $\zeta$，速率为 $\lambda>0$。成功指先到 $L$ 且发生在失效前。求成功概率及成功条件下的平均时间，并求 $\lambda\downarrow0$ 的极限。
\solution 独立失效时钟给出
\[
 u_\lambda(x)=\E_x[e^{-\lambda\tau}\ind_{\{X_\tau=L\}}].
\]
求解边值方程
\[
 \frac{\sigma^2}{2}u''-\lambda u=0,\qquad u(0)=0,\quad u(L)=1.
\]
令 $k=\sqrt{2\lambda}/\sigma$。通解由 $\sinh(kx),\cosh(kx)$ 张成，两端条件确定
\[
 \boxed{u_\lambda(x)=\frac{\sinh(kx)}{\sinh(kL)}}.
\]
为验证概率意义，对 $e^{-\lambda(t\wedge\tau)}u(X_{t\wedge\tau})$ 用伊藤公式，其漂移为0。区间上 $u,u'$ 有界，有限时刻的随机积分均值为0。布朗增量在固定时间块内足以跨越区间的概率有统一正下界，故 $\tau<\infty$ 几乎必然；支配收敛给出上述贴现终值公式。

对 $\lambda>0$，函数 $\tau e^{-\lambda\tau}$ 在非负时间上有界，在 $\lambda$ 的小邻域内也有共同上界，可在期望内求导。故
\[
 \E_x(\tau\mid\text{成功})=-\partial_\lambda\log u_\lambda(x)
 =\boxed{\frac{L\coth(kL)-x\coth(kx)}{\sigma^2k}}.
\]
利用 $\sinh z\sim z$、$\coth z=1/z+z/3+O(z^3)$，
\[
 \boxed{u_\lambda(x)\to\frac xL},\qquad
 \boxed{\E_x(\tau\mid\text{成功})\to\frac{L^2-x^2}{3\sigma^2}}.
\]
\insight 独立失效把路径按持续时间以 $e^{-\lambda\tau}$ 加权，成功路径偏向更快到达；对成功概率的对数求导恰好得到条件时间。

\problem{CR01}{泊松分流与首次到达}{经典例题}
事件按速率 $\lambda>0$ 的泊松过程到达。每次到达独立标记为 A 或 B，概率分别为 $p$ 和 $q=1-p$，其中 $0<p<1$。记两类计数为 $N_A(t),N_B(t)$。证明它们是独立泊松过程，并求给定总计数为 $n$ 时 A 类计数的分布。再令 $\tau$ 为首次 A 到达时刻，$K=N_B(\tau)$；求 $K$ 的分布及给定 $K=k$ 时 $\tau$ 的分布。
\solution 固定 $t$，先对总计数取条件期望，联合概率母函数为
\[
 \E[z^{N_A(t)}w^{N_B(t)}]
 =\exp\{\lambda t(pz+qw-1)\}
 =e^{\lambda pt(z-1)}e^{\lambda qt(w-1)}.
\]
它分解为速率 $\lambda p,\lambda q$ 的两个泊松计数母函数。原过程在不交时间区间上独立，标记也独立，因此相同分解对任意有限组不交区间同时成立，证明两个计数过程独立。给定总计数 $n$，其 $n$ 个独立标记直接给出
\[
 \boxed{N_A(t)\mid N_A(t)+N_B(t)=n\sim\operatorname{Bin}(n,p)}.
\]
设 $\alpha=\lambda p,\beta=\lambda q$。首次 A 到达满足 $\tau\sim\operatorname{Exp}(\alpha)$，且与整个 B 过程独立。因此
\[
 \Prob(K=k,\tau\in\dd t)
 =\alpha e^{-(\alpha+\beta)t}\frac{(\beta t)^k}{k!}\dd t,
 \qquad k=0,1,\ldots.
\]
利用 $\int_0^\infty t^ke^{-\lambda t}\dd t=k!/\lambda^{k+1}$，得到
\[
 \boxed{\Prob(K=k)=pq^k}.
\]
将联合密度除以此概率，得到
\[
 \boxed{f_{\tau\mid K=k}(t)=\frac{\lambda^{k+1}}{k!}t^ke^{-\lambda t},\quad t>0},
 \qquad \boxed{\E(\tau\mid K=k)=\frac{k+1}{\lambda}}.
\]
这正是 $k+1$ 个速率 $\lambda$ 的独立指数间隔之和：首次 A 之前恰有 $k$ 个 B。
\insight 固定时间内的独立计数，在按其中一类的首次到达截停后产生依赖；应先写联合密度，再作条件化。

\problem{CR02}{随机时点的剩余寿命}{经典例题}
更新过程从时刻0开始，相邻更新间隔独立，分别以概率 $1/2$ 取1和3。在 $[0,t]$ 上独立均匀选取观察时刻 $U_t$。令 $L_t$ 为包含该时刻的完整更新间隔长度，$A_t$ 为距上次更新的时间，$R_t$ 为距下次更新的时间。求 $t\to\infty$ 时三者的联合分布，并计算极限下的 $\E R$、$\Prob(R>2)$ 和 $\Cov(A,R)$。
\solution 长度为 $\ell$ 的间隔被随机时点选中的机会与其持续时间成正比。严格地，在完整的前 $n$ 个间隔中，长度为 $\ell$ 的总占用时间除以总时间，依强大数定律趋于
\[
 \frac{\ell\Prob(L=\ell)}{\E L},\qquad \E L=2.
\]
区间末端未完成部分至多长3，对时间比例的影响趋于0。因此极限长度 $L_*$ 的分布为
\[
 \boxed{\Prob(L_*=1)=\frac14,\qquad\Prob(L_*=3)=\frac34}.
\]
在长度为 $\ell$ 的间隔中，均匀观察时点的位置均匀分布，故条件于 $L_*=\ell$，
\[
 \boxed{A\sim U[0,\ell],\qquad R=\ell-A}.
\]
以上混合分布确定了三者的联合分布。原变量均有界，故分布极限也给出所需矩的极限。由于 $\E L^2=5,\E L^3=14$，
\[
 \boxed{\E R=\E A=\frac{\E L^2}{2\E L}=\frac54},\qquad
 \boxed{\Prob(R>2)=\frac{\E(L-2)^+}{\E L}=\frac14}.
\]
此外，长度为 $\ell$ 的一个间隔贡献的年龄与剩余寿命乘积积分为
$\int_0^\ell u(\ell-u)\dd u=\ell^3/6$，所以
\[
 \E(AR)=\frac{\E L^3}{6\E L}=\frac76,
 \qquad \boxed{\Cov(A,R)=\frac76-\left(\frac54\right)^2=-\frac{19}{48}}.
\]
\insight “随机选一个间隔”和“随机选一个时点所在的间隔”是不同抽样机制。后一机制先对长度加权，再在选中的间隔内均匀取位置。

\problem{CR03}{单服务台的忙期}{经典例题}
顾客按速率 $\lambda$ 的泊松过程到达，服务时间独立服从速率 $\mu$ 的指数分布，与到达过程独立，且 $0<\lambda<\mu$。一个服务台按先到先服务工作，队列从一名顾客开始。令 $B$ 为系统首次变空的时间。求 $B$ 的拉普拉斯变换、均值和方差，并说明临界情形 $\lambda=\mu$ 的均值为何不同。
\solution 将每位顾客服务期间到达的顾客看作其子代。条件于服务时间 $S$，子代数为 $\operatorname{Pois}(\lambda S)$，故平均子代数为 $\rho=\lambda/\mu<1$。第 $j$ 代顾客数的均值为 $\rho^j$，总顾客数的均值为 $1/(1-\rho)$；因其有限，总顾客数几乎必然有限。每位顾客是否存在不依赖其自身尚未开始的服务时间，由 Tonelli 定理，
\[
 \E B=\frac1\mu\sum_{j=0}^\infty\rho^j=\frac1{\mu-\lambda}.
\]
令 $b(s)=\E e^{-sB}$，$s\ge0$。首位顾客服务期间的到达顾客各自引出独立、同分布的后代服务总量。忙期长度等于树上所有服务时间之和，处理顺序不影响该总量。条件于 $S$，得到
\[
 b(s)=\E\left[e^{-sS}b(s)^{N(S)}\right]
 =\frac{\mu}{\mu+s+\lambda[1-b(s)]}.
\]
故 $\lambda b^2-(\lambda+\mu+s)b+\mu=0$。选择 $s>0$ 时落在 $(0,1)$ 内的根，并在0连续延伸，
\[
 \boxed{b(s)=\frac{\lambda+\mu+s-\sqrt{(\lambda+\mu+s)^2-4\lambda\mu}}{2\lambda}}.
\]
对 $s>0$ 求导合法，因为 $B^je^{-sB}$ 有界。对二次方程两次求导，写 $D_s=2\lambda b(s)-\lambda-\mu-s$，有
\[
 D_s b'(s)=b(s),\qquad
 D_s b''(s)=2b'(s)-2\lambda[b'(s)]^2.
\]
当 $s\downarrow0$，$b(s)\to1,D_s\to\lambda-\mu$。由非负变量的单调收敛，导数的有限极限给出
\[
 \boxed{\E B=\frac1{\mu-\lambda}},\qquad
 \E B^2=\frac{2\mu}{(\mu-\lambda)^3},
\]
\[
 \boxed{\Var B=\frac{\lambda+\mu}{(\mu-\lambda)^3}}.
\]
临界时，后代数生成函数为 $f(z)=\mu/[\mu+\lambda(1-z)]$。有限代内灭绝的概率从0开始反复迭代 $f$，极限为 $[0,1]$ 内的最小不动点；$\lambda=\mu$ 时该不动点为1。因此忙期仍几乎必然结束，但每代顾客数均值都为1，上述非负服务时间级数发散，故 $\E B=\infty$。
\insight 忙期包含服务期间新到达的全部后代。只将初始顾客的服务时间与一个平均到达数相乘，会漏掉后续各代。


\par\Needspace{8\baselineskip}
\subsection*{小节练习}
\exercise{OR25}{布朗积分的相关系数}
设 $(W_t)_{t\ge0}$ 为标准布朗运动，$T>0$ 为确定常数。定义
\[
 X=\int_0^T W_t\,dt,\qquad
 Y=\int_0^T t\,dW_t,
\]
其中第二个积分为伊藤积分。求 $X$ 与 $Y$ 的相关系数。

\exercise{OR26}{二维布朗运动的首次到轴位置}
二维标准布朗运动从 $(1,1)$ 出发，即位置为 $(1+B_t^{(1)},1+B_t^{(2)})$，其中两个分量是独立的一维标准布朗运动。首次碰到 $y$ 轴时停止。求停止位置位于 $y$ 轴正半轴的概率。


\section{练习题}
\noindent\textbf{要点与方法.}\enspace
练习中先注明计时与计费约定：是否包含初始访问、终止一步和空状态自环。再选择前缀状态、方向状态、后验参数或条件协方差，避免只按题目表面套用公式。

解出候选值后，检查边界、概率归一化及策略可行性。停时求和要辨认继续事件依赖哪些过去信息；条件时间和条件分布则须使用相应的加权，不能沿用无条件结果。
\par\smallskip



\exercise{RS1}{自重叠子串的等待}
独立抛掷公平硬币，从空记录开始，首次出现连续子串 HTHT 时停止。用前缀状态求平均抛掷次数，并说明 HTH 后再出现 H 时应保留哪个状态。
\exercise{RS2}{清空队列的尾概率}
队列每步独立以概率 $1/4$ 增加1，以概率 $1/2$ 服务一项（空时不动），以概率 $1/4$ 清空。求平稳时队长至少为2的概率，以及从空队列开始在正时间首次返回空队列的均值。
\exercise{RS3}{隔步观察的不可辨识性}
二态对称马尔可夫链每步以未知概率 $\alpha\in(0,1)$ 切换状态。精确知道它的两步切换概率为 $3/8$。求所有可能的 $\alpha$，并判断只记录偶数时刻的状态是否能区分它们。
\exercise{RS4}{递归任务数量的方差}
A 任务以概率 $1/2$ 产生两个 B，否则不产生子任务；B 以概率 $1/3$ 产生一个 A，否则不产生子任务。各任务选择独立。从一个 A 开始，求处理任务总数的均值、二阶矩与方差，并论证矩有限。

\exercise{RH1}{重启次数与总费用}
整数过程从1出发，到0即止。每步以概率 $1/2$ 重置到1，否则等概率左右各走一步。重置收费3，移动收费1。求平均停止步数、平均重置次数与平均总费用。
\exercise{RH2}{三级回撤与最高点}
公平整数游走从0出发，相对历史最高点首次回撤3格时停止。求时间均值与方差，以及停止前最高点至少为4的概率。
\exercise{RH3}{初始方向的信息}
在 $\{0,\ldots,6\}$ 上到端点即止的持续游走，每步以概率 $3/4$ 保持前一方向，否则反向。从3出发，分别在前一方向为右、为左时求先到6的概率。若前一方向本身等概率取左右且未被观察，成功概率是多少？
\exercise{RH4}{首次重复的费用}
三种类型独立等概率抽取，费用分别为1、2、4，首次出现重复类型时停止，最后一次也收费。求总费用均值，以及停止时已见不同类型数的均值。

\exercise{RC1}{禁止卖空的看跌支付}
股票现价1，一期后为2或 $1/2$，现金利率为0，允许分数持仓。期末支付为 $(1-S_1)^+$。求允许借款和卖空时的复制价格与持仓；若股票持仓必须非负，求所有路径均覆盖该支付的最小初始资金。
\exercise{RC2}{两次机会的学习价值}
未知硬币概率先验为 $U[0,1]$，给定概率后各次独立。H 收益 $+1$，T 收益 $-1$，已得收益保留，最多抛两次且随时可停。求最优规则和初始价值。
\exercise{RC3}{相关报价的替换}
免费得到 $X\sim U[0,1]$。可以直接接受，也可支付 $1/8$ 将奖值替换为 $1-X$，且只允许一次替换、不能召回。求最优规则与期望净收益。替换值与原值的关系已知。
\exercise{RC4}{两种判定目标的测试顺序}
三项独立测试的通过概率为 $(1/2,3/4,1/4)$，成本为 $(1,2,1)$。分别在“判定全部通过”和“判定至少一项通过”的目标下求最优顺序与最小期望成本；一旦结论确定即停止。

\pagebreak
\exercise{RG1}{指数积分的二次预测}
对标准布朗运动令 $I=\int_0^1 e^{W_t-t/2}\dd t$。在 $\alpha+\beta W_1+\gamma(W_1^2-1)$ 中求均方误差最小的预测及其误差。
\exercise{RG2}{双条件下的负协方差}
令 $A=\int_0^2W_t\dd t$。给定标准布朗运动满足 $W_2=0$ 且 $A=1$，求 $W_1$ 的条件分布及 $W_{1/2},W_{3/2}$ 的条件协方差。
\exercise{RG3}{桥的最大偏离}
标准布朗运动给定 $W_2=0$。求 $\max_{0\le t\le2}W_t\ge1$ 的条件概率及该最大值的条件均值。
\exercise{RG4}{失效前的成功首达}
标准布朗运动从1出发，到0或2即止；独立失效时间服从参数 $1/2$ 的指数分布。求失效前先到2的概率与成功条件下的平均时间，并与没有失效时钟时比较。
